\documentclass[12pt,letterpaper,onecolumn]{exam}
\usepackage{amsmath,amsfonts,amssymb,amsthm}
\usepackage[lmargin=71pt,tmargin=1.2in]{geometry}

\begin{document}
\thispagestyle{empty}
\begingroup
    \centering
    \LARGE Machine Learning\\
    \LARGE Homework 3: Reference Solutions\\[0.5em]
    \large Nov. 2026\\[0.5em]
    \normalsize Released: Nov. 24, 2026 \qquad Due: Dec. 16, 2026\par
\endgroup
\noindent\rule{\textwidth}{0.4pt}\par
\printanswers
\renewcommand{\solutiontitle}{\noindent\textbf{Ans:}\enspace}

\begin{questions}
    \question Regularized backpropagation.
    \begin{solution}
        Use the notation of Problem 1 in the assignment.
        The regularizer depends directly on the weights; it has no
        dependence on $z^1$, $h$, or $z^2$ when these are treated as
        intermediate variables in the computation graph.
        \begin{parts}
            \part Write $p=\hat y$. For $j,k\in\{1,\ldots,q\}$,
            \[
                \frac{\partial p_k}{\partial z^2_j}
                =p_k(\mathbf{1}_{\{k=j\}}-p_j),
                \qquad
                \frac{\partial L_{\mathrm{CE}}}{\partial p_k}
                =-\frac{y_k}{p_k}.
            \]
            By the chain rule and $\sum_k y_k=1$,
            \[
                \frac{\partial L_{\mathrm{CE}}}{\partial z^2_j}
                =-\sum_{k=1}^q y_k(\mathbf{1}_{\{k=j\}}-p_j)
                =p_j-y_j.
            \]
            Thus
            \[
                \boxed{\delta^2=\hat y-y\in\mathbb{R}^q.}
            \]

            \part Since $z^2_k=\sum_{i=1}^n W^2_{ik}h_i+b^2_k$,
            \[
                \frac{\partial L_{\mathrm{CE}}}{\partial W^2_{ik}}
                =h_i\delta^2_k,\qquad
                \frac{\partial L_{\mathrm{CE}}}{\partial b^2_k}
                =\delta^2_k.
            \]
            Also,
            $\partial(\lambda\|W^\ell\|_F^2/2)/\partial W^\ell
            =\lambda W^\ell$ for $\ell=1,2$. Therefore
            \[
                \boxed{\frac{\partial J}{\partial W^2}
                =h(\delta^2)^T+\lambda W^2
                \in\mathbb{R}^{n\times q},\qquad
                \frac{\partial J}{\partial b^2}
                =\delta^2\in\mathbb{R}^q.}
            \]
            Propagating through the output linear map gives
            \[
                \frac{\partial L_{\mathrm{CE}}}{\partial h}
                =W^2\delta^2\in\mathbb{R}^n.
            \]
            The sigmoid acts element-wise, so
            \[
                \boxed{\delta^1
                =(W^2\delta^2)\odot h\odot(\mathbf{1}-h)
                \in\mathbb{R}^n,}
            \]
            where $\mathbf{1}\in\mathbb{R}^n$ is the all-ones vector.
            Finally, $z^1_i=\sum_{r=1}^d W^1_{ri}x_r+b^1_i$ implies
            \[
                \frac{\partial L_{\mathrm{CE}}}{\partial W^1_{ri}}
                =x_r\delta^1_i,\qquad
                \frac{\partial L_{\mathrm{CE}}}{\partial b^1_i}
                =\delta^1_i.
            \]
            Adding the direct regularization gradient gives
            \[
                \boxed{\frac{\partial J}{\partial W^1}
                =x(\delta^1)^T+\lambda W^1
                \in\mathbb{R}^{d\times n},\qquad
                \frac{\partial J}{\partial b^1}
                =\delta^1\in\mathbb{R}^n.}
            \]
            The biases have no additional regularization term.

            \part Using the current parameters to evaluate both errors,
            \begin{align*}
                W^1_{\mathrm{new}}
                &=(1-\eta\lambda)W^1-\eta x(\delta^1)^T,\\
                W^2_{\mathrm{new}}
                &=(1-\eta\lambda)W^2-\eta h(\delta^2)^T.
            \end{align*}
            Compared with the unregularized updates, each weight matrix
            receives the additional term $-\eta\lambda W^\ell$.
            For $0<\eta\lambda<1$, the factor $1-\eta\lambda$
            shrinks the current weights before the data-gradient step.
            There is no such decay term for the biases.
        \end{parts}
    \end{solution}

    \newpage
    \question One gradient descent step with ReLU and $\lambda=0$.
    \begin{solution}
        All numerical results below are rounded to four decimal places;
        the updates use unrounded intermediate values.
        \begin{parts}
            \part The hidden pre-activation and activation are
            \[
                z^1=
                \begin{bmatrix}0.4&0.2\\-0.3&0.5\end{bmatrix}
                \begin{bmatrix}1\\-1\end{bmatrix}
                +\begin{bmatrix}0.1\\0.2\end{bmatrix}
                =\begin{bmatrix}0.3\\-0.6\end{bmatrix},
                \qquad
                h=\begin{bmatrix}0.3\\0\end{bmatrix}.
            \]
            Thus
            \[
                z^2=
                \begin{bmatrix}0.6&-0.2\\-0.4&0.3\end{bmatrix}
                \begin{bmatrix}0.3\\0\end{bmatrix}
                +\begin{bmatrix}0.1\\-0.1\end{bmatrix}
                =\begin{bmatrix}0.28\\-0.22\end{bmatrix}.
            \]
            The prediction and loss are
            \[
                \hat y=\frac{1}{e^{0.28}+e^{-0.22}}
                \begin{bmatrix}e^{0.28}\\e^{-0.22}\end{bmatrix}
                \approx\begin{bmatrix}0.6225\\0.3775\end{bmatrix},
            \]
            \[
                L_{\mathrm{CE}}=-\log\hat y_2
                =\log(1+e^{0.5})\approx 0.9741.
            \]

            \part Let $s=1/(1+e^{-0.5})\approx0.6224593312$.
            The output error is
            \[
                \delta^2=\hat y-y=
                \begin{bmatrix}s\\-s\end{bmatrix}
                \approx\begin{bmatrix}0.6225\\-0.6225\end{bmatrix}.
            \]
            Since both hidden pre-activations are nonzero, the ReLU
            derivative is unambiguous:
            \[
                \operatorname{ReLU}'(z^1)
                =\begin{bmatrix}1\\0\end{bmatrix},\qquad
                W^2\delta^2=\begin{bmatrix}s\\-0.5s\end{bmatrix}.
            \]
            Hence
            \[
                \delta^1=(W^2\delta^2)\odot
                \operatorname{ReLU}'(z^1)
                =\begin{bmatrix}s\\0\end{bmatrix}
                \approx\begin{bmatrix}0.6225\\0\end{bmatrix}.
            \]
            The four gradients are
            \begin{align*}
                \frac{\partial L_{\mathrm{CE}}}{\partial W^1}
                &=x(\delta^1)^T
                =\begin{bmatrix}s&0\\-s&0\end{bmatrix}
                \approx\begin{bmatrix}0.6225&0\\-0.6225&0\end{bmatrix},\\
                \frac{\partial L_{\mathrm{CE}}}{\partial b^1}
                &=\delta^1
                \approx\begin{bmatrix}0.6225\\0\end{bmatrix},\\
                \frac{\partial L_{\mathrm{CE}}}{\partial W^2}
                &=h(\delta^2)^T
                =\begin{bmatrix}0.3s&-0.3s\\0&0\end{bmatrix}
                \approx\begin{bmatrix}0.1867&-0.1867\\0&0\end{bmatrix},\\
                \frac{\partial L_{\mathrm{CE}}}{\partial b^2}
                &=\delta^2
                \approx\begin{bmatrix}0.6225\\-0.6225\end{bmatrix}.
            \end{align*}

            \part Applying $\theta_{\mathrm{new}}
            =\theta-0.1\,\partial L_{\mathrm{CE}}/\partial\theta$
            to each parameter, using the same initial gradients,
            \begin{align*}
                W^1_{\mathrm{new}}
                &=\begin{bmatrix}0.4-0.1s&-0.3\\0.2+0.1s&0.5\end{bmatrix}
                \approx\begin{bmatrix}0.3378&-0.3000\\0.2622&0.5000\end{bmatrix},\\
                b^1_{\mathrm{new}}
                &=\begin{bmatrix}0.1-0.1s\\0.2\end{bmatrix}
                \approx\begin{bmatrix}0.0378\\0.2000\end{bmatrix},\\
                W^2_{\mathrm{new}}
                &=\begin{bmatrix}0.6-0.03s&-0.4+0.03s\\-0.2&0.3\end{bmatrix}
                \approx\begin{bmatrix}0.5813&-0.3813\\-0.2000&0.3000\end{bmatrix},\\
                b^2_{\mathrm{new}}
                &=\begin{bmatrix}0.1-0.1s\\-0.1+0.1s\end{bmatrix}
                \approx\begin{bmatrix}0.0378\\-0.0378\end{bmatrix}.
            \end{align*}

            \part Hidden neuron 2 is inactive because $z^1_2=-0.6<0$.
            Its activation and ReLU derivative are both zero.
            Therefore $\delta^1_2=0$: column 2 of
            $\partial L_{\mathrm{CE}}/\partial W^1$ and entry 2 of
            $\partial L_{\mathrm{CE}}/\partial b^1$ are zero.
            Also, $h_2=0$, so row 2 of
            $\partial L_{\mathrm{CE}}/\partial W^2$ is zero.
            The corresponding parameters are unchanged in this step.
            Both output-bias gradients are nonzero. These conclusions
            concern this sample and these initial parameters; they do
            not imply that neuron 2 is inactive for every input.
        \end{parts}
    \end{solution}
\end{questions}
\end{document}
